\section{Method}

\subsection{Problem Formulation}

Let a frozen RGB Gaussian prior be
$\Gprior=\{(\boldsymbol\mu_j,\boldsymbol\Sigma_j,
\alpha_j,\mathbf{c}_j)\}_{j=1}^{N_G}$ and let the mapping set
$\Queries=\{(I_q,K_q,T_q)\}$ contain RGB images with calibrated poses. We
reconstruct a localization map
\begin{equation}
  \Anchors=\{(\mathbf{x}_i,\mathbf{d}_i,\ell_i)\}_{i=1}^{N_A},
\end{equation}
where $\mathbf{x}_i\in\R^3$ is PnP geometry, $\mathbf{d}_i$ is one deployment
descriptor, and $\ell_i$ stores Gaussian support lineage. In general,
\begin{equation}
  \text{Gaussian ID}\neq\text{track ID}\neq\text{anchor ID}.
\end{equation}
The prior remains frozen throughout localization-map reconstruction.

\subsection{KCS/GWFF Bootstrap}

Native SuperPoint observations are extracted from every mapping image. We
select a broad Gaussian-supported scaffold with keypoint-consensus sampling
(KCS): a primitive must receive support from multiple views, viewpoint bins,
and trajectory bins. Geometry-weighted feature fusion (GWFF) initializes each
selected descriptor from visible observations, weighting incidence and view
geometry and trimming low-consensus descriptor observations. This bootstrap is
an initialization, not the final landmark identity.

\subsection{Track-First Evidence and Geometry}

We build reciprocal, cycle-consistent cross-view associations from the native
observations. A track $t\in\Tracks$ is accepted only when its observations span
sufficient views and viewpoint bins. Its 3D position is estimated independently
of Gaussian centers by robust multi-view triangulation:
\begin{equation}
  \mathbf{x}_t^*=
  \arg\min_{\mathbf{x}}\sum_{(q,\mathbf{u})\in t}
  \rho\!\left(\left\|\pi(K_q,T_q,\mathbf{x})-\mathbf{u}\right\|_2\right),
\end{equation}
subject to parallax, reprojection, and covariance gates. Gaussian rasterization
then returns visible primitive IDs and composition mass at each observation.
Aggregating this provenance gives surface support $\ell_t$ without replacing
$\mathbf{x}_t^*$ by a primitive center.

\subsection{Self-Localization-Guided Reconstruction}

At iteration $s$, the current map retrieves a top-$K$ candidate set for each
native mapping keypoint. Known mapping poses label outcomes relative to the
current map. The reconstruction objective separates:
\begin{equation}
  \mathcal{L}_{\mathrm{SL}} =
  \lambda_k\mathcal{L}_{\mathrm{keep}}+
  \lambda_s\mathcal{L}_{\mathrm{swap}}+
  \lambda_m\mathcal{L}_{\mathrm{miss}}+
  \lambda_a\mathcal{L}_{\mathrm{attractor}}+
  \lambda_l\mathcal{L}_{\mathrm{local}}+
  \lambda_t\mathcal{L}_{\mathrm{trust}}.
\end{equation}
The keep term protects already correct high-margin assignments. Swap and miss
terms rank a geometrically valid positive above the selected wrong candidate or
retrieve a positive absent from the current top ranks. The global-attractor
term suppresses landmarks that collect false matches across queries. A weak
local peak objective sharpens spatially nearby alternatives, while a protected
set prevents such updates from damaging pose-critical correspondences. Because
candidates are regenerated from the current descriptors, this is an online
matching loop during offline reconstruction rather than static post-processing.

\subsection{Localization Topology Distillation}

The candidate universe contains track-centric anchors and Gaussian-supported
coverage candidates. We first select a reliable track core using cross-view and
geometric evidence. A coverage reserve then adds candidates only where mapping
queries lack enough usable rows, with pose-information scoring preventing a
purely appearance-driven map. The resulting topology has a scene-dependent
size determined by the frozen selection rule rather than a manually chosen
per-scene target.

On this fixed compact topology, a short low-rank metric update uses the complete
positive teacher and group-robust weighting. The deployed descriptor remains a
single vector per anchor; no family prototypes, learned sampler, dense stage,
or test-time rendering is required.

\subsection{Sparse Deployment}

For a query image, native SuperPoint produces keypoints and descriptors. A
shared metric transforms query and map descriptors, cosine top-1 matching
produces one correspondence per keypoint, and one standard PoseLib PnP/RANSAC
call estimates the pose. All correspondences are passed to the solver without a
landmark-side match cap or a learned test-time selector.
